\documentclass[12pt]{article}

\usepackage[utf8]{inputenc}
\usepackage{CJKutf8}
\usepackage{amsmath,amssymb}
\usepackage{amsthm}
\usepackage{geometry}
\geometry{
  a4paper,
  top=25mm,
  bottom=25mm,
  left=20mm,
  right=20mm
}
\usepackage{xcolor}
\usepackage{url}
\usepackage[unicode,colorlinks=true,linkcolor=blue,urlcolor=blue]{hyperref}
\usepackage{xurl}
\Urlmuskip=0mu plus 2mu
\sloppy
\usepackage{tocloft}

% 目录中 section 条目使用点线填充到页码
\renewcommand{\cftsecleader}{\cftdotfill{\cftdotsep}}

% 基本定理环境
\newtheorem{definition}{定义}[section]
\newtheorem{axiom}{公理}[section]
\newtheorem{assumption}{假设}[section]
\newtheorem{theorem}{定理}[section]
\newtheorem{proposition}{命题}[section]
\newtheorem{corollary}{推论}[section]
\newtheorem{lemma}{引理}[section]
\newtheorem{remark}{备注}[section]
\newtheorem{Certificate}{证书}[section]

% 记号（仅用于本笔记自洽引用，避免依赖主稿宏定义）
\newcommand{\GFramework}{\textsf{G-Framework}}
\newcommand{\Layer}{\ensuremath{\ell}}
\newcommand{\Jac}{\operatorname{Jac}}
\newcommand{\Ob}{\ensuremath{\mathsf{Ob}}}
\newcommand{\JacGap}{\ensuremath{\mathsf{JacGap}}}
\newcommand{\RR}{\mathbb{R}}
\newcommand{\NN}{\mathbb{N}}
\newcommand{\code}[1]{\path{#1}}

% 避免少量不可控的换行溢出（尤其是混排 CJK 与等宽字体时）
\setlength{\emergencystretch}{6em}

\begin{document}
\begin{CJK*}{UTF8}{gkai}

\title{边缘算子封闭性、比安基基准与变分生成的同伦修复塔\\（工作笔记：形式系统版）}
\author{Gao Zheng（高政）}
\date{2026-01-11}
\maketitle

\section*{前言}
\addcontentsline{toc}{section}{前言}

本文的目标不是单独说明边缘算子、比安基恒等式或变分生成塔，而是在 \GFramework{}
的修复基础中，把低阶边缘算子、雅可比失配、上同调障碍、同伦塔、Bianchi/Jacobi
同一性与变分生成组织成一个闭合判定系统。若这些对象分散出现，边缘算子容易只被
理解为局部微分，变分塔容易只被理解为迭代技巧；本文将它们统一为“什么样的边缘
结构足以支撑同伦闭合”的问题。

本文的第一层立场是边缘算子的充分性判定。低阶边缘算子如果无法消除三元障碍类，
就会表现为雅可比失配或上同调失败。引入新边缘算子并不只是增加工具，而是为了
恢复闭合、吸收障碍，并使系统能够进入更高阶同伦塔。

本文的第二层立场是 Bianchi 与 Jacobi 的统一。交换子恒等式与曲率恒等式在本文中
被视为同一闭合需求的不同投影：前者偏代数一致性，后者偏几何曲率约束。通过这一
统一，边缘算子的失败可以被同时理解为代数失配与几何失闭合。

本文的第三层立场是变分塔的生成作用。修复后边缘算子的结构不是一次性指定，而是
由高阶一致性逐层确定。变分塔给出生成这些高阶边缘修复项的机制，使修复不再只是
补丁，而是有方向、有层级、有闭合证书的递归过程。

本文的第四层立场是低阶失败的可上升化。低阶 Jacobi 失败并不意味着系统整体不可
用，而意味着当前层级缺少足够的边缘算子或高阶同伦数据。本文将这种失败上升为
可修复的障碍类，使其进入更高阶同伦塔，而不是直接判定结构崩溃。

本文的第五层立场是 exactness 的修复语义。当 Jacobiator 可以成为某个边缘算子的
边界时，它在上同调中被平凡化，低阶严格失败就被高阶同伦吸收。这个过程给出了
“失败仍可修复”的精确条件，也为后续 JacGap 作为证书型残差提供基础语义。

本文的第六层立场是变分生成的可计算性。高阶边缘算子若只能被抽象声明，就很难服务
工程系统；本文强调它们可以通过变分泛函、生成规则或递归塔结构被构造。由此，修复
塔不仅是存在性语言，也是可计算修复方向的来源。

本文的第七层立场是为 \(L_\infty\) 修复塔提供前置骨架。后续更高阶 \(L_\infty\) 算子作为
同伦修复塔时，需要知道低阶失败如何进入高阶、边缘算子如何被新增、Bianchi/Jacobi
如何被统一解释。本文正提供这些基础接口。

因此，本文在整体 \GFramework{} 中承担的是“边缘算子闭合与变分修复塔”的基础锚定。
它向前连接同伦代数与 JacGap，向中间连接上同调障碍、Bianchi/Jacobi 与高阶一致性，
向后为 \(L_\infty\) 修复塔、语义边缘算子和统计命中中的雅可比间隙提供闭合语义。概括地说，
本文把“边缘算子是否够用”推进为“能否通过同伦闭合与变分生成塔吸收障碍”的形式系统。

更具体地说，本文把“闭合失败”从错误消息转化为结构诊断。低阶边缘算子不够用时，
系统不会直接停机，而是得到一个可上升的障碍类；这个障碍类可以请求新边缘算子、
新变分项或新高阶同伦数据。这样，失败本身变成修复塔的输入。

从方法论上看，本文是 G-Framework 的边界生成稿。边缘算子定义什么是边界，变分塔
定义如何生成新的边界修复项，Bianchi/Jacobi 统一说明这些边界为何需要闭合。后续
所有讨论 JacGap、上同调障碍和高阶修复的稿件，都在不同程度上调用了这个边界语义。

从整体谱系看，本文把抽象代数闭合推向可计算修复方向。它既不只停留在“存在高阶
结构”，也不直接跳到工程算法，而是说明高阶结构如何由变分生成、如何被边缘算子
组织、如何在同伦塔中逐层消化障碍。

\setcounter{tocdepth}{3}
\setcounter{secnumdepth}{3}
\tableofcontents
\clearpage

%%%%%%%%%%%%%%%%%%%%%%%%%%%%%%%%%%%%%%%%%%%%%%%%%%%%%%%%%%%%
\section{概览：把“边缘算子是否够用”翻译为同伦代数命题链}
%%%%%%%%%%%%%%%%%%%%%%%%%%%%%%%%%%%%%%%%%%%%%%%%%%%%%%%%%%%%

\begin{remark}[本章定位]
本章把一条工程直觉命题链，翻译为同伦代数的标准语言，并写成可检验的证明结构：
\begin{enumerate}
  \item \textbf{低阶边缘算子不封闭}
        $\Longleftrightarrow$
        二元括号的 Jacobi 恒等式失效，出现 Jacobiator（雅可比缺陷）；
  \item \textbf{引入新的边缘算子后恢复上同调封闭性}
        $\Longleftrightarrow$
        Jacobiator 成为 $l_1$-边界（exact），因此在上同调里 Jacobi 严格成立；
  \item \textbf{“比安基恒等式 $\Rightarrow$ 雅可比恒等式”}
        $\Longleftrightarrow$
        在“联络—曲率”的改写下，交换子 Jacobi 与 Bianchi 是同一恒等式的两种表达；
  \item \textbf{同伦塔揭示修复后的结构}
        $\Longleftrightarrow$
        更高阶算子 $(l_3,l_4,\dots)$ 逐层吸收阻碍，形成 $L_\infty$（强同伦李）结构塔；
  \item \textbf{变分法构成边缘算子的形态}
        $\Longleftrightarrow$
        以作用量（尤其 BV/BRST 语境）为生成元，主方程给出 $Q^2=0$，从而“自动”产出满足同伦恒等式的 $(l_n)$。
\end{enumerate}
与本仓库中更系统的同伦修复塔论证可对照：
\code{NoteBook/notebook1/tex12_pub/gframework_linfty_homotopy_repair_formal_system.tex}。\cite{HomotopyRepairTex12}
\end{remark}

%%%%%%%%%%%%%%%%%%%%%%%%%%%%%%%%%%%%%%%%%%%%%%%%%%%%%%%%%%%%
\section{低阶边缘算子：雅可比失配与三元障碍类}
%%%%%%%%%%%%%%%%%%%%%%%%%%%%%%%%%%%%%%%%%%%%%%%%%%%%%%%%%%%%

\subsection{设定：$(l_1,l_2)$、导子性与 Jacobiator}

\begin{definition}[分级空间与低阶边缘算子]
设 $V=\bigoplus_{n\in\mathbb{Z}}V^n$ 为 $\mathbb{Z}$-分级向量空间。
对齐次元素 $x\in V$，记其次数为 $|x|$。

给定两个边缘算子（低阶结构算子）：
\begin{enumerate}
  \item $l_1:V\to V$，次数为 $+1$；
  \item $l_2:\wedge^2 V\to V$，次数为 $0$（分级反对称二元括号）。
\end{enumerate}
本章默认 $l_1^2=0$，从而上同调
\[
  H(V,l_1):=\ker(l_1)/\mathrm{im}(l_1)
  \quad\text{可定义。}
\]
\end{definition}

\begin{assumption}[$l_1$ 对 $l_2$ 的导子性（兼容性）]
\label{ass:l1-derivation-l2}
对任意齐次 $x,y\in V$，有分级导子公式
\begin{equation}
\label{eq:l1-derivation}
  l_1\bigl(l_2(x,y)\bigr)
  =
  l_2(l_1x,y)+(-1)^{|x|}\,l_2(x,l_1y).
\end{equation}
\end{assumption}

\begin{definition}[Jacobiator（雅可比缺陷）]
给定 $l_2$，定义其 Jacobiator 为
\begin{equation}
\label{eq:jacobiator-def}
\begin{aligned}
  J(x,y,z)
  &:=
  l_2(l_2(x,y),z)
  +(-1)^{|x|(|y|+|z|)}\,l_2(l_2(y,z),x)
  \\
  &\qquad
  +(-1)^{|z|(|x|+|y|)}\,l_2(l_2(z,x),y),
\end{aligned}
\end{equation}
其中 $x,y,z$ 均为齐次。
若 $J\equiv 0$，则称 $l_2$ 严格满足 Jacobi 恒等式（链级严格封闭）；否则称其不封闭。
\end{definition}

\subsection{命题：Jacobiator 是 $l_1$-闭的（比安基基准）}

\begin{proposition}[Bianchi-闭性：$l_1(J)=0$]
\label{prop:bianchi-closed-jacobiator}
在 $l_1^2=0$ 与 Assumption~\ref{ass:l1-derivation-l2} 成立时，对任意齐次 $x,y,z\in V$，有
\[
  l_1\bigl(J(x,y,z)\bigr)=0.
\]
\end{proposition}

\begin{proof}
将 $l_1$ 作用到 \eqref{eq:jacobiator-def} 的三项上，并对每一项反复使用导子公式 \eqref{eq:l1-derivation} 展开。
展开后出现两类项：
\begin{enumerate}
  \item $l_2(\,\cdot\,,l_1(\cdot))$ 型项：它们在三项循环求和与分级符号约定下两两抵消；
  \item $l_2(l_1l_2(\cdot,\cdot),\cdot)$ 型项：继续用 \eqref{eq:l1-derivation} 展开后，剩余项包含 $l_1^2$，由 $l_1^2=0$ 消去。
\end{enumerate}
因此 $l_1(J)=0$。
\end{proof}

\begin{definition}[三元障碍类（上同调中的 Jacobi 失配类）]
\label{def:jacobiator-class}
若 $x,y,z\in\ker(l_1)$，则由 Proposition~\ref{prop:bianchi-closed-jacobiator} 得 $J(x,y,z)\in\ker(l_1)$，
因此可定义其上同调类
\[
  [J](x,y,z):=[J(x,y,z)]\in H(V,l_1).
\]
称 $[J]$ 为低阶边缘算子 $(l_1,l_2)$ 的\textbf{三元障碍类}：它刻画“Jacobi 失效是否能在上同调层面消去”。
\end{definition}

\begin{corollary}[拓扑连通口径：相关上同调为零 $\Rightarrow$ 障碍类平凡]
\label{cor:connectedness-kills-obstruction}
若在相关次数上 $H(V,l_1)=0$（例如可理解为“无三阶上同调障碍/拓扑连通的占位口径”），
则对任意 $l_1$-闭的 $x,y,z$ 都有 $[J](x,y,z)=0$。
\end{corollary}

\begin{proof}
由 Definition~\ref{def:jacobiator-class}，$[J](x,y,z)\in H(V,l_1)$。
若该同调群为零，则任何同调类皆为零，结论成立。
\end{proof}

%%%%%%%%%%%%%%%%%%%%%%%%%%%%%%%%%%%%%%%%%%%%%%%%%%%%%%%%%%%%
\section{引入新边缘算子：上同调封闭性与障碍修复}
%%%%%%%%%%%%%%%%%%%%%%%%%%%%%%%%%%%%%%%%%%%%%%%%%%%%%%%%%%%%

\subsection{三元修复算子 $l_3$：把 Jacobiator 变成边界}

\begin{definition}[三元边缘算子（修复算子）]
给定一个三元边缘算子
\[
  l_3:\wedge^3V\to V,
  \qquad
  \deg(l_3)=-1.
\]
称 $l_3$ 为 $l_2$ 的\textbf{雅可比修复算子}，若对任意 $l_1$-闭的齐次 $x,y,z\in\ker(l_1)$，有
\begin{equation}
\label{eq:arity3-repair}
  J(x,y,z)=l_1\bigl(l_3(x,y,z)\bigr).
\end{equation}
\end{definition}

\begin{theorem}[障碍可消 $\Longleftrightarrow$ 存在 $l_3$ 完成修复（上同调封闭性）]
\label{thm:obstruction-vanish-iff-l3}
在 Proposition~\ref{prop:bianchi-closed-jacobiator} 的前提下，下列两者等价：
\begin{enumerate}
  \item 对任意 $l_1$-闭三元组 $(x,y,z)$，有 $[J](x,y,z)=0$；
  \item 存在三元边缘算子 $l_3$ 使得修复方程 \eqref{eq:arity3-repair} 对任意 $l_1$-闭 $x,y,z$ 成立。
\end{enumerate}
\end{theorem}

\begin{proof}
$(2)\Rightarrow(1)$：若 \eqref{eq:arity3-repair} 成立，则 $J(x,y,z)$ 为 $l_1$-边界，故其同调类为零。

$(1)\Rightarrow(2)$：固定任意 $l_1$-闭的齐次 $x,y,z$。
由 $[J](x,y,z)=0$，存在某个元素 $h(x,y,z)\in V$ 使得 $J(x,y,z)=l_1(h(x,y,z))$。
令 $l_3(x,y,z):=h(x,y,z)$（对齐次数约定），则 \eqref{eq:arity3-repair} 成立。
对所有 $l_1$-闭三元组重复该构造即可得到所需 $l_3$。
\end{proof}

\subsection{推论：修复后上同调中的 Jacobi 恒等式严格成立}

\begin{corollary}[上同调封闭：$H(V,l_1)$ 上诱导出严格李括号]
\label{cor:cohomology-lie-bracket}
在 Assumption~\ref{ass:l1-derivation-l2} 与 Theorem~\ref{thm:obstruction-vanish-iff-l3} 的条件下，
二元括号 $l_2$ 在上同调上诱导出良定义括号
\[
  [\![x]\!]\,,[\![y]\!]\ \longmapsto\ [\![l_2(x,y)]\!],
\]
并且该诱导括号满足严格 Jacobi 恒等式。
\end{corollary}

\begin{proof}
良定义性来自导子性 \eqref{eq:l1-derivation}：改变代表元只会改变一个 $l_1$-边界。
Jacobi 恒等式在上同调上的成立来自 Theorem~\ref{thm:obstruction-vanish-iff-l3}：
对 $l_1$-闭输入，Jacobiator $J$ 为 $l_1$-边界，故其同调类为零。
\end{proof}

%%%%%%%%%%%%%%%%%%%%%%%%%%%%%%%%%%%%%%%%%%%%%%%%%%%%%%%%%%%%
\section{同伦塔：修复后边缘算子的结构由高阶一致性逐层确定}
%%%%%%%%%%%%%%%%%%%%%%%%%%%%%%%%%%%%%%%%%%%%%%%%%%%%%%%%%%%%

\subsection{从 $(l_1,l_2,l_3)$ 到 $(l_4,l_5,\dots)$：逐层吸收更高阶阻碍}

\begin{remark}[从“一个 $l_3$”到“一个塔”]
Theorem~\ref{thm:obstruction-vanish-iff-l3} 只处理了最低阶（$3$ 元）Jacobi 缺陷。
若要求全体系在同伦意义下完全封闭，则需要更高阶算子 $(l_4,l_5,\dots)$ 逐层吸收更高阶阻碍。
这正是 $L_\infty$ 结构的“同伦塔”语义：每一层 $l_{n+1}$ 都对应于“上一层连贯性方程的剩余缺陷”。
\end{remark}

\begin{definition}[$L_\infty$ 的塔式读法（接口化）]
称一族多线性映射 $\{l_n\}_{n\ge 1}$（其中 $l_1$ 次数 $+1$，$l_n$ 次数 $2-n$）
构成一个 $L_\infty$ 结构，若其满足全部高阶连贯性方程（可写成余导子 $Q$ 的条件 $Q^2=0$）。
在该读法下：
\begin{itemize}
  \item $l_1$ 给出边界算子；
  \item $l_2$ 给出二元边缘算子（括号）；
  \item $l_3$ 吸收 $l_2$ 的 Jacobi 缺陷；
  \item 更高 $l_n$ 吸收更高阶阻碍项 $\Ob_n$。
\end{itemize}
\end{definition}

\begin{proposition}[逐层延拓判据：下一层阻碍为边界当且仅当可引入 $l_{N+1}$]
\label{prop:extension-obstruction-edge}
固定一个截断的同伦代数数据 $(V,l_1,\{l_k\}_{2\le k\le N})$，并假设其满足到阶 $N$ 的一致性方程。
记 $N{+}1$ 阶一致性方程中\emph{不含} $l_{N+1}$ 的剩余项为 $\Ob_{N+1}$。
则在 $l_1$-闭输入上，存在 $l_{N+1}$ 使得 $N{+}1$ 阶一致性成立的充要条件是
\[
  \Ob_{N+1}=-\,l_1(l_{N+1})\qquad(\text{在 }l_1\text{-闭输入上}).
\]
因此，“同伦塔揭示结构”可被理解为：$l_{N+1}$ 的形态由 $\Ob_{N+1}$ 的边界原像决定。
\end{proposition}

\begin{proof}
这是 $L_\infty$ 连贯性方程的标准“阻碍 = 边界”读法：
将 $N{+}1$ 阶恒等式整理为“含 $l_{N+1}$ 的项 + 不含 $l_{N+1}$ 的项 = 0”，
并在 $l_1$-闭输入上消去“某个输入先被 $l_1$ 作用再进入 $l_{N+1}$”的项，即得该判据。
更系统的推导可对照 \code{NoteBook/notebook1/tex12_pub/gframework_linfty_homotopy_repair_formal_system.tex}
中的同名延拓命题。\cite{HomotopyRepairTex12}
\end{proof}

%%%%%%%%%%%%%%%%%%%%%%%%%%%%%%%%%%%%%%%%%%%%%%%%%%%%%%%%%%%%
\section{Bianchi 与 Jacobi：交换子恒等式与曲率恒等式的同一性}
%%%%%%%%%%%%%%%%%%%%%%%%%%%%%%%%%%%%%%%%%%%%%%%%%%%%%%%%%%%%

\begin{lemma}[交换子 Jacobi $\Rightarrow$ Bianchi（经典几何型表述）]
\label{lem:jacobi-implies-bianchi}
设 $\{D_i\}$ 为一组算子（可理解为协变导数/联络算子），并定义曲率
\[
  F_{ij}:=[D_i,D_j].
\]
则交换子的 Jacobi 恒等式
\[
  [D_i,[D_j,D_k]]+\text{cyc}=0
\]
等价地推出 Bianchi 形式
\[
  [D_i,F_{jk}]+[D_j,F_{ki}]+[D_k,F_{ij}]=0.
\]
\end{lemma}

\begin{proof}
将 $F_{jk}=[D_j,D_k]$ 代回 Jacobi 恒等式即可得到。
\end{proof}

\begin{remark}[“基于 Bianchi 的 Jacobi”在本章的含义]
Lemma~\ref{lem:jacobi-implies-bianchi} 说明：当边缘算子被解释为“联络/协变导数”时，
“代数闭合性（Jacobi）”与“曲率一致性（Bianchi）”是同一结构的两种投影。
在同伦语境中，该对应被推广为“广义 Bianchi”（见下一节）：其内容等价于整套同伦恒等式 $Q^2=0$。
\end{remark}

%%%%%%%%%%%%%%%%%%%%%%%%%%%%%%%%%%%%%%%%%%%%%%%%%%%%%%%%%%%%
\section{广义 Bianchi：Maurer--Cartan 曲率与同伦闭合}
%%%%%%%%%%%%%%%%%%%%%%%%%%%%%%%%%%%%%%%%%%%%%%%%%%%%%%%%%%%%

\begin{definition}[Maurer--Cartan 曲率（$L_\infty$ 口径）]
给定 $L_\infty$ 数据 $(V,\{l_n\})$。
对一个元素 $A\in V$（可视为“联络/场”），定义其曲率（或缺陷）为
\[
  \mathcal{F}(A):=\sum_{n\ge 1}\frac{1}{n!}\,l_n(A,\dots,A),
\]
其零点方程 $\mathcal{F}(A)=0$ 即 Maurer--Cartan 方程。
\end{definition}

\begin{proposition}[广义 Bianchi 恒等式（接口式）]
\label{prop:generalized-bianchi}
若 $\{l_n\}$ 满足 $L_\infty$ 一致性方程（等价于 $Q^2=0$），则对任意 $A$ 有
\[
  \sum_{n\ge 0}\frac{1}{n!}\,l_{n+1}(A,\dots,A,\mathcal{F}(A))=0.
\]
\end{proposition}

\begin{proof}
该恒等式是 $Q^2=0$ 在“把 $A$ 视为同一输入反复插入”时的分量展开形式，
可视为同伦结构中的 Bianchi-闭性陈述：曲率/缺陷在同伦导数下满足循环为零的约束。
\end{proof}

%%%%%%%%%%%%%%%%%%%%%%%%%%%%%%%%%%%%%%%%%%%%%%%%%%%%%%%%%%%%
\section{变分法生成边缘算子：BV/BRST 作为形态来源（接口化）}
%%%%%%%%%%%%%%%%%%%%%%%%%%%%%%%%%%%%%%%%%%%%%%%%%%%%%%%%%%%%

\subsection{主方程 $\Rightarrow Q^2=0$：从作用量到同伦微分}

\begin{definition}[BV 括号与 BV 微分（接口化）]
设存在一个带奇 Poisson 括号（BV 括号）的分级代数 $(\mathcal{A},\{\cdot,\cdot\}_{BV})$。
给定作用量（泛函）$S\in\mathcal{A}$，定义算子
\[
  Q:=\{S,\cdot\}_{BV}.
\]
\end{definition}

\begin{definition}[BV 主方程（Classical master equation）]
称 $S$ 满足 BV 主方程，若
\[
  \{S,S\}_{BV}=0.
\]
\end{definition}

\begin{proposition}[主方程 $\Rightarrow Q^2=0$]
若 $S$ 满足 BV 主方程，则 $Q$ 满足 $Q^2=0$。
\end{proposition}

\begin{proof}
由 BV 括号的 Jacobi 恒等式，
\[
  Q^2(\cdot)=\{S,\{S,\cdot\}\}
  =\tfrac12\{\{S,S\},\cdot\},
\]
因此当 $\{S,S\}=0$ 时有 $Q^2=0$。
\end{proof}

\subsection{从 $Q$ 的展开到 $(l_n)$：边缘算子形态的“自动产出”}

\begin{theorem}[变分生成的同伦塔（形式结论）]
\label{thm:variational-generates-linfty}
在上述接口下，若 $Q=\{S,\cdot\}_{BV}$ 满足 $Q^2=0$，并且 $Q$ 在某个背景点的泰勒展开存在，
则其多线性展开系数可以被组织为一族运算 $\{l_n\}_{n\ge 1}$，
使其满足 $L_\infty$ 的全部一致性方程。
因此，“变分法构成边缘算子的形态”可被严格读作：
\[
  \{S,S\}_{BV}=0 \Longrightarrow Q^2=0 \Longrightarrow (l_1,l_2,l_3,\dots)\ \text{满足同伦 Jacobi/Bianchi。}
\]
\end{theorem}

\begin{proof}
将 $Q$ 作为一个次数 $+1$ 的生成算子，$Q^2=0$ 即给出一组二次关系。
把这些关系按输入的多重次数分解，得到逐阶一致性方程；
将每一阶的多线性分量记为 $l_n$，即可得到同伦塔。
该论证属于 BV/BRST 语境中的标准“主方程 $\Rightarrow$ 代数结构”读法；本章将其作为接口化定理使用。
\end{proof}

%%%%%%%%%%%%%%%%%%%%%%%%%%%%%%%%%%%%%%%%%%%%%%%%%%%%%%%%%%%%
\section{结语：障碍修复、拓扑连通与边缘算子结构}
%%%%%%%%%%%%%%%%%%%%%%%%%%%%%%%%%%%%%%%%%%%%%%%%%%%%%%%%%%%%

\begin{remark}[一句话封口]
在本章的形式系统中：
\begin{itemize}
  \item “不封闭”被编码为三元障碍类 $[J]\neq 0$；
  \item “引入新边缘算子后封闭”被编码为存在 $l_3$ 使 $J=l_1(l_3)$，从而上同调括号严格满足 Jacobi；
  \item “比安基基准”对应于 $l_1(J)=0$（缺陷必须是可追踪的闭缺陷）以及其在几何/曲率口径下的等价表达；
  \item “同伦塔揭示结构”对应于每阶阻碍 $\Ob_{n+1}$ 的边界原像确定了 $l_{n+1}$ 的形态；
  \item “变分法给出形态”对应于主方程 $\{S,S\}=0$ 强制 $Q^2=0$，从而把 $(l_n)$ 作为展开系数自动产出。
\end{itemize}
\end{remark}

\clearpage

\appendix

% ============================================================
% 附录A：全文标准化形式表达（Jacobi--Bianchi--同伦塔--变分生成）
% ============================================================
\section{边缘算子不封闭与上同调障碍修复：Jacobi--Bianchi--同伦塔与变分生成}
\label{sec:edge-op-cohomology-repair}

\subsection{符号约定（Shuffles 与 Koszul 符号）}
\label{subsec:edge-op-conventions}

设 $V=\bigoplus_{k\in\mathbb{Z}}V^k$ 为分次向量空间；齐次元 $x\in V$ 的次数记为 $|x|$（上同调次数约定）。

\paragraph{定义 D0（$(i,n-i)$-shuffle 集合）.}
对 $n\ge 1$ 与 $0\le i\le n$，记
\[
\mathrm{Sh}(i,n-i)
=\bigl\{\sigma\in S_n:\ \sigma(1)<\cdots<\sigma(i),\ \sigma(i+1)<\cdots<\sigma(n)\bigr\}.
\]

\paragraph{定义 D1（Koszul 符号）.}
给定齐次元 $(x_1,\dots,x_n)$ 与 $\sigma\in S_n$，定义
\[
\chi(\sigma;x_1,\dots,x_n)
\;:=\;\mathrm{sgn}(\sigma)\cdot (-1)^{\sum_{p<q,\ \sigma(p)>\sigma(q)}|x_p||x_q|}.
\]
下文简写为 $\chi(\sigma)$（默认依赖 $(x_1,\dots,x_n)$）。

\subsection{边缘算子不封闭 $\Rightarrow$ 雅可比障碍类}
\label{subsec:edge-op-jacobi-obstruction}

\paragraph{定义 D2（低阶边缘算子）.}
称一对多线性算子 $(l_1,l_2)$ 为\emph{低阶边缘算子}，若
\[
l_1:V\to V,\quad \deg(l_1)=+1;\qquad
l_2:\wedge^2 V\to V,\quad \deg(l_2)=0.
\]
并假设
\[
l_1^2=0.
\tag{E1}
\label{eq:E1-diff}
\]
于是可定义上同调 $H(V,l_1)$。

\paragraph{定义 D3（导子兼容/Leibniz）.}
称 $l_1$ 与 $l_2$ \emph{兼容}，若对任意齐次元 $x,y\in V$，
\[
l_1\,l_2(x,y)=l_2(l_1x,y)+(-1)^{|x|}l_2(x,l_1y).
\tag{E2}
\label{eq:E2-derivation}
\]
（旁注：这是 $L_\infty$ 的 $n=2$ 一致性条件。）

\paragraph{定义 D4（Jacobiator/三元障碍）.}
定义 $l_2$ 的三元缺陷（Jacobiator）为
\[
\Ob_3^{(2)}(x_1,x_2,x_3)
:=\sum_{\sigma\in \mathrm{Sh}(2,1)}\chi(\sigma)\,
l_2\!\bigl(l_2(x_{\sigma(1)},x_{\sigma(2)}),\,x_{\sigma(3)}\bigr).
\tag{E3}
\label{eq:E3-ob3}
\]
当 $\Ob_3^{(2)}\equiv 0$ 时，$l_2$ 在链层面满足严格雅可比；一般情形下 $\Ob_3^{(2)}$ 即“边缘算子不封闭”的精确刻画。

\paragraph{命题 P1（障碍是 $l_1$-闭的）.}
在 \eqref{eq:E1-diff} 与 \eqref{eq:E2-derivation} 前提下，对任意齐次元 $x_1,x_2,x_3$，
\[
l_1\Ob_3^{(2)}(x_1,x_2,x_3)
\;=\;
\sum_{r=1}^3 (-1)^{|x_1|+\cdots+|x_{r-1}|}\,
\Ob_3^{(2)}(x_1,\dots,l_1x_r,\dots,x_3).
\tag{E4}
\label{eq:E4-ob3-closed}
\]
特别地，若 $l_1x_1=l_1x_2=l_1x_3=0$，则 $l_1\Ob_3^{(2)}(x_1,x_2,x_3)=0$。

\paragraph{证明.}
对 \eqref{eq:E3-ob3} 逐项作用 $l_1$，并对每一项用 \eqref{eq:E2-derivation} 把 $l_1$ 下推到最内层输入；
再对 $\mathrm{Sh}(2,1)$ 的三项按 Koszul 符号配对整理，可得所有“$l_1$ 作用在内层括号”的项完全改写为“对输入逐个插入 $l_1$”的和，
得到 \eqref{eq:E4-ob3-closed}。证毕。

\begin{Certificate}[三元障碍类可消（$l_1$-exact）的可检验证书]
\label{cert:edge-ob3-exact}
对 $l_1$-闭的齐次元 $x_1,x_2,x_3$，$\Ob_3^{(2)}(x_1,x_2,x_3)$ 定义了上同调类
\[
\bigl[\Ob_3^{(2)}(x_1,x_2,x_3)\bigr]\in H(V,l_1).
\]
若存在次数 $-1$ 的三线性映射 $B_3:\wedge^3V\to V$ 使得对所有 $l_1$-闭输入
\[
l_1 B_3(x_1,x_2,x_3)=-\Ob_3^{(2)}(x_1,x_2,x_3),
\tag{C1}
\label{eq:C1-ob3-exact}
\]
则该障碍类为零，即 $\bigl[\Ob_3^{(2)}\bigr]=0$（在闭输入上）。
\end{Certificate}

\begin{Certificate}[收缩同伦给出障碍可消的充分条件（连通/可缩情形）]
\label{cert:edge-contraction-suffices}
若存在次数 $-1$ 的线性算子 $h:V\to V$（收缩同伦）与投影 $\pi:V\to H(V,l_1)$ 使得
\[
l_1h+hl_1=\mathrm{Id}_V-\iota\pi,
\tag{C2}
\label{eq:C2-contraction}
\]
其中 $\iota:H(V,l_1)\to V$ 为某个同调代表的嵌入，则对任意 $l_1$-闭输入
\[
\Ob_3^{(2)}(x_1,x_2,x_3)=l_1\bigl(-h\,\Ob_3^{(2)}(x_1,x_2,x_3)\bigr),
\]
从而 \eqref{eq:C1-ob3-exact} 成立（取 $B_3=-h\circ \Ob_3^{(2)}$ 的“按输入三线性化”的版本）。
\end{Certificate}

\subsection{引入新边缘算子 $l_3$：上同调障碍修复与 $L_\infty$ 的 $n=3$ 恒等式}
\label{subsec:edge-op-add-l3}

\paragraph{定义 D5（新边缘算子 $l_3$）.}
引入三元边缘算子
\[
l_3:\wedge^3V\to V,\qquad \deg(l_3)=-1.
\]

\begin{Certificate}[$n=3$ 的同伦雅可比（亦可视为广义 Bianchi）]
\label{cert:edge-linfty-n3}
称 $(l_1,l_2,l_3)$ 满足 $L_\infty$ 的 $n=3$ 一致性条件，若对任意齐次元 $x_1,x_2,x_3$，
\[
l_1l_3(x_1,x_2,x_3)
\;+\;
\sum_{\sigma\in\mathrm{Sh}(1,2)}\chi(\sigma)\,
l_3\!\bigl(l_1x_{\sigma(1)},x_{\sigma(2)},x_{\sigma(3)}\bigr)
\;+\;
\Ob_3^{(2)}(x_1,x_2,x_3)
\;=\;0.
\tag{E5}
\label{eq:E5-linfty-3}
\]
（旁注：这条式子等价于“Jacobiator 在链层面是 $l_1$-边界（并包含输入插入 $l_1$ 的修正项）”。）
\end{Certificate}

\paragraph{定理 T1（障碍为零 $\Rightarrow$ 存在 $l_3$ 完成修复）.}
在 \eqref{eq:E1-diff} 与 \eqref{eq:E2-derivation} 前提下，若 Certificate~\ref{cert:edge-ob3-exact} 成立（即闭输入上的障碍类为零），
则存在 $l_3$ 使得 Certificate~\ref{cert:edge-linfty-n3} 的恒等式 \eqref{eq:E5-linfty-3} 成立。
并且由 $l_2$ 在 $H(V,l_1)$ 上诱导的括号
\[
[a,b]_{H}:=[\,l_2(\tilde a,\tilde b)\,]\quad (\tilde a,\tilde b\ \text{为 }l_1\text{-闭代表})
\]
满足严格的分次雅可比恒等式，因此 $H(V,l_1)$ 上得到一个分次李代数结构。

\paragraph{证明.}
由 Certificate~\ref{cert:edge-ob3-exact} 取 $B_3$，令 $l_3:=B_3$（并按 \eqref{eq:E5-linfty-3} 作一致性延拓），
则 \eqref{eq:E5-linfty-3} 成立，即 Certificate~\ref{cert:edge-linfty-n3} 成立。
对 $l_1$-闭代表取上同调类，$\Ob_3^{(2)}$ 为 $l_1$-边界，从而诱导括号满足严格雅可比。证毕。

\subsection{同伦塔（Homotopy Tower）：逐层障碍与逐层修复}
\label{subsec:edge-op-homotopy-tower}

\paragraph{定义 D6（对 $n$ 线性算子的 $\delta$-边界算子）.}
对次数为 $2-n$ 的 $n$ 线性映射 $f:\wedge^nV\to V$，定义
\[
(\delta f)(x_1,\dots,x_n)
:=
l_1 f(x_1,\dots,x_n)
\;+\;
\sum_{r=1}^n (-1)^{|x_1|+\cdots+|x_{r-1}|}\,
f(x_1,\dots,l_1x_r,\dots,x_n).
\tag{E6}
\label{eq:E6-delta}
\]
（旁注：这就是“只看 $l_1$ 时”的链复形边界算子在多线性层面的提升。）

\paragraph{定义 D7（第 $n$ 阶障碍 $\Ob_n$ 的递推定义）.}
假设已构造 $l_1,\dots,l_{n-1}$，且每个 $l_k:\wedge^kV\to V$ 满足 $\deg(l_k)=2-k$。
定义第 $n$ 阶障碍（仅由低于 $n$ 阶算子拼出）为
\[
\Ob_n(x_1,\dots,x_n)
:=
\sum_{\substack{i+j=n+1\\ 2\le i,j\le n-1}}
\ \sum_{\sigma\in\mathrm{Sh}(i,n-i)}
\chi(\sigma)\,(-1)^{i(j-1)}\,
l_j\!\bigl(l_i(x_{\sigma(1)},\dots,x_{\sigma(i)}),x_{\sigma(i+1)},\dots,x_{\sigma(n)}\bigr).
\tag{E7}
\label{eq:E7-obn-rec}
\]
于是第 $n$ 阶一致性可被写成一条“线性层级方程”
\[
\delta l_n \;+\;\Ob_n \;=\;0.
\tag{E8}
\label{eq:E8-level-n}
\]

\begin{Certificate}[归纳假设：到 $n-1$ 阶的一致性]
\label{cert:edge-tower-ih}
假设对所有 $m\le n-1$，下式成立：
\[
\sum_{i+j=m+1}\ \sum_{\sigma\in\mathrm{Sh}(i,m-i)}
\chi(\sigma)\,(-1)^{i(j-1)}\,
l_j\!\bigl(l_i(x_{\sigma(1)},\dots,x_{\sigma(i)}),x_{\sigma(i+1)},\dots,x_{\sigma(m)}\bigr)=0.
\tag{IH$_{m}$}
\label{eq:IH-m}
\]
（旁注：这就是“同伦塔已封闭到 $n-1$ 层”。）
\end{Certificate}

\begin{Certificate}[第 $n$ 阶障碍是 $\delta$-闭的（从而给出障碍上同调类）]
\label{cert:edge-obn-closed}
在 Certificate~\ref{cert:edge-tower-ih} 前提下，有
\[
\delta \Ob_n = 0.
\tag{C4a}
\label{eq:C4a-obn-closed}
\]
特别地，对 $l_1$-闭输入，$\Ob_n(x_1,\dots,x_n)$ 是 $l_1$-闭元，从而定义障碍类 $[\Ob_n]\in H(V,l_1)$。
\end{Certificate}

\begin{Certificate}[第 $n$ 阶可修复条件（障碍可消：$\delta$-exact）]
\label{cert:edge-obn-exact}
若存在次数 $2-n$ 的 $n$ 线性映射 $B_n:\wedge^nV\to V$ 使得
\[
\delta B_n = -\Ob_n,
\tag{C4b}
\label{eq:C4b-obn-exact}
\]
则可取 $l_n:=B_n$ 使得层级方程 \eqref{eq:E8-level-n} 成立，从而第 $n$ 阶一致性被修复。
\end{Certificate}

\paragraph{定理 T2（逐层修复：从 $n-1$ 到 $n$ 的纯调用证书）.}
在 Certificate~\ref{cert:edge-tower-ih} 前提下，若 Certificate~\ref{cert:edge-obn-exact} 成立，
则存在 $l_n$ 使得第 $n$ 阶一致性 \eqref{eq:E8-level-n} 成立；从而同伦塔从 $(l_1,\dots,l_{n-1})$ 延拓到 $(l_1,\dots,l_n)$。

\paragraph{证明.}
由 Certificate~\ref{cert:edge-obn-exact} 取 $l_n:=B_n$，立即得到 \eqref{eq:E8-level-n}。证毕。

\paragraph{归纳模板（把“修复”写成两条可复用规则）.}
对每个 $n\ge 3$，重复使用如下两条规则即可：
\[
\textbf{(R1)}\ \ \text{由 Certificate~\ref{cert:edge-tower-ih} 得 }\delta\Ob_n=0\ \ (\text{Certificate~\ref{cert:edge-obn-closed}}),
\]
\[
\textbf{(R2)}\ \ \text{若能给出 }\delta B_n=-\Ob_n\ \ (\text{Certificate~\ref{cert:edge-obn-exact}})\ \text{，则取 }l_n:=B_n\
  \text{完成修复（Theorem T2）}.
\]

\subsection{Jacobi 与 Bianchi：同一一致性的两种表述}
\label{subsec:edge-op-jacobi-bianchi}

\paragraph{定义 D8（算子视角：联络与曲率）.}
设 $\{D_i\}$ 为一族算子，定义曲率
\[
F_{ij}:=[D_i,D_j].
\]

\paragraph{引理 L1（Jacobi $\Rightarrow$ Bianchi）.}
交换子满足雅可比恒等式
\[
[D_i,[D_j,D_k]]+\text{cyc}(i,j,k)=0
\]
当且仅当曲率满足 Bianchi 形式
\[
[D_i,F_{jk}] + [D_j,F_{ki}] + [D_k,F_{ij}] = 0.
\tag{E9}
\label{eq:E9-bianchi}
\]

\paragraph{证明.}
将 $F_{jk}=[D_j,D_k]$ 代入雅可比恒等式即可得到 \eqref{eq:E9-bianchi}；反向同理。证毕。

\paragraph{说明（$L_\infty$ 语境下的“广义 Bianchi”）.}
式 \eqref{eq:E5-linfty-3} 可视为 \eqref{eq:E9-bianchi} 的同伦推广：
当 $l_2$ 的严格雅可比不成立时，$l_3$ 给出“以边界方式成立”的广义 Bianchi/广义雅可比，
从而保证上同调层面严格封闭。

\subsection{变分法生成边缘算子：主方程 $\Rightarrow$ $Q^2=0$ $\Rightarrow$ 同伦塔}
\label{subsec:edge-op-variational}

\begin{Certificate}[BV 主方程作为可检验封闭证书]
\label{cert:edge-bv-master}
在 BV 形式中，若存在作用量/泛函 $S$ 与 BV 括号 $\{\cdot,\cdot\}$ 使得
\[
\{S,S\}=0,
\tag{C5}
\label{eq:C5-master}
\]
则定义次数 $+1$ 的导子
\[
Q:=\{S,\cdot\}
\]
满足 $Q^2=0$。
\end{Certificate}

\paragraph{定理 T3（“变分法构成边缘算子的形态”）.}
在 Certificate~\ref{cert:edge-bv-master} 前提下，将 $Q$ 在某个背景点做多阶展开，其 Taylor 系数给出一族多线性算子 $\{l_n\}_{n\ge 1}$
（次数满足 $\deg(l_n)=2-n$），并自动满足全部 $L_\infty$ 一致性条件（可用 Theorem T2 的归纳模板逐层读出）。
因此边缘算子族的具体形态可由变分主方程统一生成，而同伦塔的封闭性由 $Q^2=0$ 直接保证。

\paragraph{证明.}
由 \eqref{eq:C5-master} 得 $Q^2=\{S,\{S,\cdot\}\}=\tfrac12\{\{S,S\},\cdot\}=0$。
将 $Q$ 视为余导子/高阶导出括号的生成元，其展开系数满足的全部二次约束等价于 $Q^2=0$ 的分量方程，
从而得到 $L_\infty$ 的一切恒等式（亦可按 Theorem T2 的归纳模板逐层读出）。证毕。

\clearpage

\renewcommand{\refname}{\texorpdfstring{参考文献}{References}}
\begin{thebibliography}{99}

\bibitem{HomotopyRepairTex12}
\GFramework 工作笔记：
\emph{在 G-Framework 中：更高阶 $L_\infty$ 算子作为同伦修复塔}
（\code{NoteBook/notebook1/tex12_pub/gframework_linfty_homotopy_repair_formal_system.tex}；比安基基准、JacGap 证书封闭与同伦塔等）。

\end{thebibliography}

\end{CJK*}
\end{document}
